\section{A real quartic class certificate for the conductor-four Tate residues}
\label{sec:conductor4-quartic-class-certificate}

Put $p=269$, $K=\mathbb Q(\sqrt{-3})$, $t=\sqrt[3]p>0$ and
\[
 \begin{gathered}
 s=\sqrt p>0,\qquad w=\sqrt{(p+13s)/2}>0,\qquad\\
 F=\mathbb Q(w),\quad L=F(t),\quad M=KL.
 \end{gathered}
 \tag{G152C.1}
\]
Let $\epsilon=\chi_{-4}$ be the primitive odd quadratic character
of conductor four, $\psi=\omega_p^{67}$ with $i\equiv82\pmod p$,
and $\xi=\epsilon\psi$.
Then $\psi(2)=-i$, $\xi$ is primitive even of conductor $4p=1076$,
and $\xi^2=\tau_p$ is the real quadratic character of conductor $p$.
Characters are zero at integers not coprime to their primitive conductor.

\begin{theorem}[AC-G152-class: the real conductor-four quartic class sector]
\label{AC-G152-class}
For the degree-twelve field in (G152C.1),
\[
                         \operatorname{Cl}(L)[269]=0.
 \tag{G152C.2}
\]
The nontrivial cubic sector of $\operatorname{Cl}(M)[269]$ is
therefore zero.
Seven explicitly certified units generate a subgroup of finite
index $m$ in $\mathcal O_L^\times/\mu_L$ and satisfy
\[
                             h_L=7m,\qquad269\nmid m.
 \tag{G152C.3}
\]
Neither $m=1$, $h_L=7$, a complete class group, nor vanishing
of the trivial cubic class sector of $M$ is asserted.
The original-coefficient strict and actual-window statements are separate.
\end{theorem}

\begin{proof}
\proofstep{1. The exact quartic character and its two generator orientations.}
Use the already exactly verified Gaussian periods from AC-G147-class.
For $z=\zeta_p$, set
\[
 \eta_j=\sum_{a\in2^j(\mathbb F_p^\times)^4}z^a,
 \qquad v=\eta_0-\eta_2,\quad v'=\eta_1-\eta_3.
 \tag{G152C.4}
\]
Their complete integer-polynomial certificate, retained with its
hash, proves
\[
 2v^2=-p-13s,\qquad vv'=-5s,\qquad
 v^4+pv^2+25p=0,
 \tag{G152C.5}
\]
where the finite quadratic Gauss sign gives $s=\sqrt p$.
The generator $g_p:z\mapsto z^2$ sends
$v\mapsto v'\mapsto-v\mapsto-v'$ and has $\psi(g_p)=-i$.
The quartic period field has conductor $p$ and is disjoint
from $\mathbb Q(i)$, since the latter ramifies at two.
In their product $C_2\times C_4$, let $c_i$ denote conjugation
on $\mathbb Q(i)$.  The kernel of $\epsilon\psi$ is the
order-two subgroup $\langle(c_i,g_p^2)\rangle$, generated by
that single ordered pair.
It fixes $iv$, whose square is $(p+13s)/2$.
The norm of this square is $25p$, a rational nonsquare;
thus $iv$ has degree four and generates precisely the field of $\xi$.
Its possible sign relative to $w>0$ does not change the odd
automorphism formulas below.
Consequently
\[
 \begin{gathered}
 P(X)=X^4-269X^2+6725,\\
 \qquad F=\mathbb Q(\zeta_{4p})^{\ker\xi},\qquad
 D(F)=(4p)^2p=2^4p^3.
 \end{gathered}
 \tag{G152C.6}
\]
Here the final equality is the abelian conductor--discriminant formula.
Primitivity is literal: the factor at two has conductor four,
and the factor at $p$ has order four on tame inertia.
Both are nontrivial and their conductors are coprime.
The square character is $\tau_p$, so the unique quadratic
subfield is $\mathbb Q(\sqrt p)$.

Let $b_-$ fix $i$ and act as $g_p$ on the period field.
Since $w\,b_-(w)=5s$ and $s=(2w^2-p)/13$, one obtains
\[
 \begin{aligned}
 b_-(w)&=(w^3-219w)/65,&\xi(b_-)&=-i,\\
 b_+(w):=b_-^{-1}(w)&=(219w-w^3)/65,&\xi(b_+)&=i.
 \end{aligned}
 \tag{G152C.7}
\]
Both squares negate $w$.
The exact dictionary verifies these polynomial identities modulo $P$;
the corresponding CRT representative for $b_-$ is $809$ modulo $1076$,
congruent to one modulo four and two modulo $p$.
The four roots of $P$ are
$\pm\sqrt{(p+13s)/2}$ and $\pm\sqrt{(p-13s)/2}$,
all real since $p>13s$.
The two inverse characters $\xi,\xi^{-1}$ define this same field,
with the orientations in (G152C.7) kept explicit.

\proofstep{2. The degree-twelve field and its ramified local data.}
The splitting field $K(t)$ has group $S_3$ and unique quadratic
subfield $K$, different from $\mathbb Q(\sqrt p)$.
It is disjoint from $F$, so
\[
 \operatorname{Gal}(M/\mathbb Q)=S_3\times C_4,
 \quad [L:\mathbb Q]=12,
 \quad(r_1(L),r_2(L))=(4,4),\quad |\mu_L|=2.
 \tag{G152C.8}
\]
The signature combines the four real places of $F$ with
the one real and two complex embeddings of $\mathbb Q(t)$.
A real embedding excludes all roots of unity except $\pm1$.

The cubic field has discriminant $-3p^2$.
At three, the Hensel equation $z+3z^2+3z^3=-30$
gives the root $-(1+3z)$ of $X^3-p$ in $\mathbb Q_3$;
the remaining quadratic factor is $K_3$, of discriminant exponent one.
At $p$ the cubic field is tame Eisenstein, and its polynomial
discriminant excludes other primes besides three and $p$.
In particular it is unramified at two.
The quartic discriminant exponent at two is four, so unramified
base change by the cubic gives exponent twelve in $L$.
At three the quartic is unramified, giving exponent four.
The polynomial $P$ is Eisenstein at $p$; combining its tame
ramification degree four with the cubic degree three gives
$e=12,f=1$, with discriminant exponent eleven.
Hence
\[
 D:=D(L)=2^{12}3^4p^{11}
       =177059831617390347261447643090944.
 \tag{G152C.9}
\]
There is one prime of $L$ above $p$.
Since $K_p/\mathbb Q_p$ is unramified quadratic, $M$ has
one prime above $p$, with $e=12,f=2$ and total local degree twenty-four.
These assertions do not use the provisional discriminant from extraction.

\proofstep{3. Seven exact global units and a prime-to-$269$ index.}
The unit file
\path{all_units_exact_certificate.json}
has SHA256
\begin{quote}\ttfamily\small
11b77b9b3adeeaa9290aeab824f8c48ea\\
423c0c6a8fe18d718ef8e2f2d01ca9d.
\end{quote}
It specifies seven rational vectors in the basis $t^a w^b$,
$0\le a<3$, $0\le b<4$, indexed by $4a+b$, and all their
monic integral degree-twelve relations with constant term $\pm1$.
The common denominators, in order, are
\[
                         13,13,13,3497,3497,52455,52455.
 \tag{G152C.10}
\]
They include powers of $p$ and are not discarded.
The maximum numerator has twenty-four decimal digits.

The exact regular multiplication matrices satisfy
\[
                     T^3=pI,\quad W^4-pW^2+25pI=0,
                     \quad TW=WT.
 \tag{G152C.11}
\]
The raw generator is $w-t$.
The verifier checks its resultant, irreducibility and both inverse
rational maps for $t,w$, evaluates every raw candidate unit at $W-T$,
and checks that its first column is the specified tensor vector.
It computes the exact characteristic polynomial and also checks
the same relation by rational Horner evaluation.
Monic integrality and constant term $\pm1$ prove that each value
and its inverse is integral.  Thus these are seven actual global
units without any complete-unit assertion from the extraction.

Use $\ell=206593$, $\ell-1=768p$, the primitive root $g=10$
and its order-$p$ power $h=g^{768}=49873$.
Trial division and exact power tests verify these assertions.
The three simple roots of $X^3-p$ and four simple roots of $P$
give twelve degree-one embeddings into $\mathbb Q_\ell$.
All displayed denominators are invertible modulo $\ell$.
For each embedding and unit define $c_j$ by
$\bar U_j^{768}=h^{c_j}$.
Among the twelve saved residue-logarithm rows the selected minor is
\[
 C=\begin{pmatrix}
70&221&148&264&90&227&265\\
199&121&221&90&264&4&227\\
199&148&48&90&264&265&42\\
70&48&121&264&90&42&4\\
70&221&148&61&153&194&117\\
199&121&221&153&61&152&194\\
199&148&48&153&61&117&75
\end{pmatrix}.
 \tag{G152C.12}
\]
Its exact determinant is
$-21113247082378584\equiv238\pmod p$.
These residue maps kill $p$th powers and roots of unity.
The seven units are therefore independent modulo $p$ in the
rank-seven free unit lattice.
A primitive integral relation would contradict this rank;
the unit theorem gives a full-rank subgroup, and the same
mod-$p$ invertibility gives its finite index $m$ prime to $p$.

\proofstep{4. The weighted regulator and rational isolation bound.}
Write $A=\sqrt{(p+13\sqrt p)/2}$ and
$B=\sqrt{(p-13\sqrt p)/2}$, and $\rho=p^{1/3}$.
The real embeddings are $(t,w)=(\rho,w)$ for $w=A,-A,B,-B$;
the complex-place representatives are $(\rho\omega,w)$
with those same four values.
Give real places weight one and complex places weight two,
form the eight-by-seven logarithm matrix, delete its last row
and take the absolute determinant $R_U$.
The projected weighted-log convention gives $R_U=mR_L$ exactly.
The $3000$-bit outward-ball evaluation of the rational arrays
has zero in every weighted column sum and proves
\[
                        374842116142<R_U<374842116143.
 \tag{G152C.13}
\]
No provisional regulator supplies an error bound.

For $n=12,r_2=4$, define the integer bounds
\[
 U=\left\lceil\frac{12!4^4(\lfloor\sqrt D\rfloor+1)}
                              {12^{12}3^4}\right\rceil,
 \qquad H=4^{12}\left\lceil\sqrt[3]{U^4}\right\rceil.
 \tag{G152C.14}
\]
Minkowski gives a representative of each class of norm at most
$U$, because $\pi>3$.
Euler-factor comparison and $\zeta(4/3)\le4$ give
$h_L\le U^{4/3}\zeta_L(4/3)\le4^{12}U^{4/3}\le H$.
Exact integer arithmetic gives
\[
                  U=2259316018036,
           \qquad H=497381015211158961913856.
 \tag{G152C.15}
\]
If $r=h_L/m>1$, then $m<h_L\le H$.
For an integer $c$, the inequality $|r-c|<1/H$ forces
$h_L=cm$, since a nonzero numerator $h_L-cm$ would give
$|r-c|\ge1/m>1/H$.
This avoids any lower bound for the full regulator or assumption $m=1$.

\proofstep{5. Exact Artin factors and finite even Dirichlet sums.}
Let $\varrho$ be the standard representation of the splitting
field of $X^3-p$ and let $\delta=\tau_p$.
The permutation representation for $L$ is
$(1\oplus\varrho)\otimes(1\oplus\delta\oplus\xi\oplus\bar\xi)$.
Artin formalism retains every inertia-invariant factor and gives
\[
 Z:=\operatorname*{res}_{s=1}\zeta_L(s)
 =L(1,\delta)|L(1,\xi)|^2
      L(1,\varrho)L(1,\varrho\delta)|L(1,\varrho\xi)|^2.
 \tag{G152C.16}
\]
For a primitive even character $\theta$ of conductor $q$, put
\[
 S(\theta)=\sum_{a=1}^{q-1}\bar\theta(a)
                                 \log(2\sin(\pi a/q)).
 \tag{G152C.17}
\]
The finite Fourier and Abel-limit proof in AC-G147-class,
Step4, applies to these exact primitive characters: expand
$-\log|1-r\zeta_q^a|$ in the absolutely convergent cosine
series, multiply by $\bar\theta(a)$, sum finitely, then let
$r\uparrow1$.  The finite Gauss identity and Parseval give
$L(1,\theta)=-S(\theta)/\mathfrak g(\bar\theta)$ and
$|\mathfrak g(\bar\theta)|^2=q$.
Consequently
\[
                 |L(1,\xi)|^2=|S(\xi)|^2/(4p).
 \tag{G152C.18}
\]
For $\delta$, the real Gauss sum has absolute value $\sqrt p$.
The complete finite cosine sum is a strictly positive ball,
between sixteen and seventeen, and hence has sign $+\sqrt p$.
Thus $L(1,\delta)=-S(\delta)/\sqrt p$, also strictly positive
in the calculation.
No quadratic-field class number or unit is imported.
Every finite logarithm is evaluated at a strictly positive sine ball.

\proofstep{6. The actual Hecke conductors and primitive theta formula.}
The multiplicative coefficients of $L(s,\varrho)=\sum a_n n^{-s}$
have $a_{3^v}=1$, $a_{p^v}=0$ for $v\ge1$; away from $3p$,
\[
 a_\ell=\#\{x\bmod\ell:x^3=p\}-1,\qquad
 a_{\ell^{v+1}}=a_\ell a_{\ell^v}-(-3/\ell)a_{\ell^{v-1}}.
 \tag{G152C.19}
\]
The exact cubic-residue test for $\ell\equiv1\pmod3$ gives
$a_\ell=2$ or $-1$; for $\ell\equiv2\pmod3$ the higher
coefficients alternate between one in even degree and zero in odd degree.
The twisted sequences are $a_n\delta(n)$ and $a_n\xi(n)$.
At $p$ their cubic inertia multiplied by quadratic or quartic
inertia has order six or twelve and no invariant vector.
At two the $\xi$ twist has no invariants, while the $\delta$
twist is unramified.  The sequences therefore include exactly
the bad Euler factors, with no omitted conductor correction.

Over $K$, $\varrho=\operatorname{Ind}_K^{\mathbb Q}\nu$
for a primitive cubic character of conductor $(p)$.
It is locally trivial at three by the cube-root calculation above.
The norm pullback of $\delta$ has tame order two at $p$,
so $\nu\delta$ has conductor $(p)$ again.
For the $\xi$ twist, the norm pullback at two has exact
conductor four: in the unramified quadratic extension $K_2$,
\[
 N(1+2a)\equiv1+2\operatorname{Tr}_{\mathbb F_4/\mathbb F_2}
                                    (\bar a)\pmod4.
 \tag{G152C.20}
\]
The trace is onto and $\epsilon(3)=-1$.
At $p$ the residue norm is onto, and the order-four pullback
times the order-three cubic character has tame order twelve.
All twists are unramified at three.
Thus their primitive Hecke moduli and rational conductors are
\[
 f=p,p,4p,\qquad Q=3f^2=217083,217083,3473328.
 \tag{G152C.21}
\]
In particular there is no extra factor $p$ in either twisted modulus.

The prime-power finite Fourier/Poisson proof of AC-G147-class,
Step6, applies literally to each modulus in (G152C.21).
Its hypotheses just verified are: a primitive character on
$\mathcal O_K/f$, zero on nonunits and trivial on the six
global units; $h(K)=1$ from $2\sqrt3/\pi<2$; and $3\nmid f$.
At a nonunit frequency modulo a prime power, primitivity
supplies a unit in the last principal-unit layer that fixes
the additive frequency and acts nontrivially on the character.
The sum vanishes; the unit frequencies and Parseval give Gauss
absolute value $f$.  This includes the new factor $2^2$.
The inverse-different pairing and the Euclidean dual
$2/(f\sqrt{-3})\mathcal O_K$ have the cancelling factor two
proved there, with no invertibility of two modulo $f$ presumed.
The Gaussian scalar is $1/(ft)$ and its dual exponent is
$a|\beta|^2/t$ for $a=2\pi/(\sqrt3f)$.
Conjugating the dual summation variable gives
\[
 \Theta(t)=\varepsilon t^{-1}\overline{\Theta(1/t)},
 \quad |\varepsilon|=1,\qquad
 \Theta(t)=\sum_{n\ge1}b_ne^{-ant}.
 \tag{G152C.22}
\]
Both constant terms vanish.  Mellin integration and splitting
at one give the exact identity
\[
 L(1)=\sum_{n\ge1}\frac{b_ne^{-an}}n
       +\varepsilon a\sum_{n\ge1}\bar b_n E_1(an),
 \qquad E_1(x)=\int_1^\infty e^{-xt}\frac{dt}{t}.
 \tag{G152C.23}
\]
All three theta balls at one exclude zero.
For the two real sequences the phase is one; for the complex
sequence it is the actual quotient
$\Theta(1)/\overline{\Theta(1)}$.
No root-number convention is guessed.

\proofstep{7. Outward tails, rational isolation and reflected descent.}
The recurrence implies $|b_n|\le d(n)\le n$.
For $z=e^{-a}$ and $A_N=e^{-a(N+1)}$, the omitted tails of
the first sum, second sum before multiplication by $a$, and theta
have the explicit bounds
\[
 \frac{A_N}{1-z},\qquad \frac{A_N}{a(1-z)},\qquad
                   \frac{A_N((N+1)-Nz)}{(1-z)^2}.
 \tag{G152C.24}
\]
The second follows from $E_1(x)\le e^{-x}/x$.
The $256$-bit calculation adds these full radii to each real
and, when applicable, imaginary component, with cutoffs
$12000,12000,40000$.
It retains the $3000$-bit regulator ball and all finite
Dirichlet logarithm balls.

The analytic class-number formula for four real places,
four complex places and two roots of unity gives
\[
               h_LR_L=\frac{2\sqrt D}{2^4(2\pi)^4}Z.
 \tag{G152C.25}
\]
The saved inclusion calculation proves
\[
 r=\frac{h_LR_L}{R_U}=\frac{h_L}{m}>1,
 \qquad |r-7|<1.3\cdot10^{-52},\qquad
                         1/H>2.01\cdot10^{-24}.
 \tag{G152C.26}
\]
The strict test uses the full balls, not these rounded bounds.
The integer seven is only a final selector; no provisional
class group or full unit-basis claim enters any preceding quantity.
Step4 proves $h_L=7m$ exactly, and Step3 proves $p\nmid m$.
This proves (G152C.2).

For the reflection $s_0$ fixing $L$ inside $M$, norm and
extension on class groups have composites two and $1+s_0$.
Since the $p$-primary class group of $L$ vanishes, $s_0$ acts
by $-1$ on that of $M$.
Every conjugate reflection does likewise; their products
generate the cubic rotations, which therefore act trivially.
Only the nontrivial cubic sector is concluded zero.
\end{proof}

\paragraph{Original labels and retained boundaries.}
For the actual finite label $\chi^{-j}\chi_{-4}$, $j=1,2$,
the residual Kummer quotient has quartic factor $\xi^{-1}$
at effective residue $66$ and $\xi$ at effective residue $200$.
These inverse characters define the same field $F$ above.
They are residual splitting factors only: the new character
$\psi$ is not inserted into the old finite host-character list.
The coefficient field and its original integer ring are not
changed by this class certificate.
A subsequent strict theorem must still establish both local
unit images on the actual one-place $e=12,f=2$ completion,
retaining all rational denominators and the chosen $b$ orientation,
then lift on the original coefficient DVR.
The actual-window comparison costs remain a separate requirement.

\paragraph{Sources and computation scope.}
The exact period identities, finite Gauss sign and complete
prime-power theta proof are the pinned AC-G147-class inputs;
their precise applications and changed conductors are checked above.
No redundant fresh upstream period or Poisson proof is claimed.
The analytic class-number formula and projected weighted-log
normalization retain Sutherland, MIT18.785 Lecture19 (2019),
pp.1 and4,
\url{https://math.mit.edu/classes/18.785/2019fa/LectureNotes19.pdf}.
Minkowski, Dirichlet units, Artin formalism and the abelian
conductor--discriminant formula retain their earlier audited scopes.
The new scripts verify the exact field dictionary, all seven
rational unit relations, all twelve residue maps and the strict
analytic isolation using documented python-flint0.9.0 outward balls.
The Arb implementation, including the order-one exponential
integral, is retained rather than independently re-proved.
Candidate-only PARI extraction supplies rational units and a
selector; no GRH, complete-unit assertion or \texttt{bnfcertify}
is used.
